% ===================================================================
%  第 9 节 其他联系（Other connections）
%  原文：TASI10.tex 第 2455--2472 行（PDF p.68）
% ===================================================================
\section{其他联系}
\label{sec:other_physics}

\begin{keytakeaways}
  \item 包裹 Calabi–Yau 三流形除子的 M5 膜给出具有 (0,4) 世界面超对称的有效弦，即 MSW 弦；这是又一个产生“模对象—有限群”联系的物理设置。
  \item MSW 弦的广义椭圆亏格是斜全纯 Jacobi 形式的实例，而这类模对象在推广 Thompson 月光的月光族中扮演主角。
  \item 包裹 $\PP^2$ 的两个 M5 膜给出与类数模拟模形式 \eq{hurwitz_mock} 生成函数密切相关的 mock 斜全纯形式，为 Mathieu 月光的几何理解提供了可能起点。
  \item 包裹 del Pezzo 8 与 half-$K3$ 的单个 M5 膜，其广义椭圆亏格由权 $3/2$ 模形式 $f^{(1)}= E_4\eta^{-5}$ 支配；该形式是 Mathieu 群 $M_{12}$ 超模的分次超维数，其几何与弦论解释尚待给出。
\end{keytakeaways}

\myparagraph{MSW 弦}

另一个给出与有限群相联系的模对象实例的物理设置，是包裹 Calabi–Yau 三流形除子的 M5 膜 \cite{Cheng:2017dlj}。
它们产生有效弦，即所谓的 MSW 弦，具有 (0,4) 世界面超对称 \cite{Maldacena:1997de}。

% [Tutor Note: MSW 弦指 Maldacena–Strominger–Witten 构造：M5 膜包裹 Calabi–Yau 三流形的一个除子（实的 4-循环），在剩余两个非紧方向看来是一条有效弦，世界面理论具有 (0,4) 超对称；其低能有效理论用于黑洞微观态计数。原文的 \cite{Maldacena:1997de} 即这一构造的原始文献。]

\cite{Cheng:2017dlj} 指出，MSW 弦理论的广义椭圆亏格是斜全纯 Jacobi 形式的例子
（参见 \S\ref{subsec:shjacobiforms}）；这类模对象在推广 Thompson 月光的月光族中扮演主角，
但此前其在物理中的作用一直未受强调。
特别地，\cite{Cheng:2017dlj} 考察了包裹曲面 $\PP^2$、del Pezzo 8 与 half-$K3$ 的一个或两个 M5 膜的广义椭圆亏格。

% [Tutor Note: 斜全纯 Jacobi 形式（skew-holomorphic Jacobi form）的模变换因子用 $\bar\tau$ 而非 $\tau$，这是它与通常弱 Jacobi 形式的关键差别，也是它天然出现在权 $3/2$ 月光中的原因；不要把它与 mock Jacobi 形式的“非全纯完备化”混为一谈，后者是在全纯部分之外加上影子的积分。定义见 \S\ref{subsec:shjacobiforms}。]

% [Tutor Note: half-$K3$ 不是 $K3$ 曲面本身，而指有理椭圆曲面（$\PP^2$ 在 9 点爆破，又称 dP9），可看作 $K3$ 的“一半”；阅读时若把 half-$K3$ 当成 $K3$，会错误估计 M5 膜世界体积的模形式权重。原文未给定义，细节见 \cite{Cheng:2017dlj}。]

% [Tutor Note: 作者用 “it was known” 带过而省略了推导：这里实质上是把 M5 膜世界体积上的 BPS 态计数化为某个模形式（广义椭圆亏格），再证明所得对象满足斜全纯 Jacobi 形式的模变换；完整计算见 \cite{Cheng:2017dlj}。读者若想核验“与类数生成函数密切相关”一句，应比较其 Fourier 系数与 \eq{hurwitz_mock} 的系数，而不是只看模变换性质。]

在第一种情形中，已知两个包裹曲面 $\PP^2$ 的 M5 膜导致模拟斜全纯形式（mock skew-holomorphic forms），
它们与类数模拟模形式 \eq{hurwitz_mock} 的生成函数密切相关。
这一联系暗示，$\PP^2$ 上的 M5 膜可能提供一个起点，
从那里我们可以追求对 \S\ref{ch:moonshineweight32} 所描述的（重标度的）Hurwitz 类数的 Mathieu 月光的几何理解。

% [Tutor Note: 注意区分“一个”与“两个” M5 膜：两个包裹 $\PP^2$ 的 M5 膜给出 mock 斜全纯形式，指向类数模拟模形式；单个包裹 del Pezzo 8 或 half-$K3$ 的 M5 膜则给出权 $3/2$ 模形式 $f^{(1)}=E_4\eta^{-5}$。把两种膜数对应的结果混淆，是阅读本段最常见的错误。]

对其余情形，人们发现，支配单个包裹 del Pezzo 8 与 half-$K3$ 曲面的 M5 膜的广义椭圆亏格的
权重 $3/2$ 模形式是 $f^{(1)}= E_4\eta^{-5}$。
（$E_4$ 与 $\eta$ 的定义见 \S\ref{subsec:modular_forms}。）
% [Tutor Note: 作者在此只说结果是“已知的”，未说明 $M_{12}$ 作用从何而来，也未说明为何 del Pezzo 8 与 half-$K3$ 两种不同曲面由同一个 $f^{(1)}$ 支配——这恰恰是下句“若能从几何与弦论上理解”所指的未解问题。注意不要把这里的 $M_{12}$ 与第 5 节 Mathieu 月光的 $M_{24}$ 混为一谈。]

\cite{Cheng:2017dlj} 表明，$f^{(1)}$ 是 Mathieu 群 $M_{12}$ 的一个超模的分次超维数（graded superdimension），
并且对所有 $g\in M_{12}$ 明确给出了相应的分次特征标 $f^{(1)}_g$。
若能从几何与弦论上理解这个 $M_{12}$ 超模，那将非常有趣。

% [Tutor Note: 权重核算极易出错：$E_4$ 是权 $4$ 的 Eisenstein 级数，$\eta^{-5}$ 的权为 $-5/2$，故 $f^{(1)}$ 的总权为 $4-5/2=3/2$，正好与“权 $3/2$ 模形式”相符；把 $\eta$ 的 $1/2$ 权漏掉就会得到权 $4$ 或权 $9/2$ 的错误结论。]

% [Tutor Note: 分次超维数就是超迹 $\mathrm{Str}(1)$：偶分量维数减奇分量维数。因此 $f^{(1)}$ 的 Fourier 系数是整数，可以取负值；对应的分次特征标 $f^{(1)}_g$ 是带插入 $g$ 的超迹。这与第 6 节强调的“权 $3/2$ 必须使用超模”完全一致，不要按普通迹去解读这些系数。]
